\chapter{What the validation found}\label{ch:val-findings}

Everything here was produced by the checks of the three preceding chapters on the
build under test, and every number can be re-derived from the evidence files of
Appendix~\ref{app:val-trace}. Findings are numbered D-1 upwards. Only the first has
been corrected in the page: three of the others would change stored values, and
stored values are not changed without agreement.

\section{D-1 --- a console view failed to build (corrected)}
\vhead{What happened}
The builder of the load view referred to a name that was not defined in its scope.
The reference is on the path taken whenever at least one file has been read, so the
function raised on every session that had data --- which is every real session. The
other seven views built normally, so the failure was silent: the first view of the
console simply produced nothing.

\vhead{Evidence, before the correction}

\begin{lstlisting}[language=,numbers=left,firstnumber=1,basicstyle=\ttfamily\scriptsize\color{white}]
the console views, with the corpus loaded:
  load        raised: bad is not defined
  sop         20 items
  results     4 items
  panel       54 items
  graphs      26 items
  review      13 items
  export      12 items
  integrity   1 item
\end{lstlisting}

\vhead{The correction}
The condition the line was written for is now computed immediately above it:

\begin{lstlisting}[language=,numbers=left,firstnumber=1,basicstyle=\ttfamily\scriptsize\color{white}]
const bad=RUNS.some(r=>r.acqError||(r.integrity&&r.integrity.ok===false));
\end{lstlisting}

It raises the view when a file could not be read or its checksum does not match,
and leaves it quiet otherwise. It changes no stored value and no export. All eight
views now build; the counts are in Chapter~\ref{ch:val-oracles}.

\section{Findings against the exchange-file specification}
The reader was re-checked against the published schema and format notes of the
exchange format. Each finding was measured, either on the corpus or on a file
written to isolate one clause.

\subsection*{D-2 --- three control types are not recognised}
The specification defines eight sample types. The page maps six, one of which the
specification does not define.

{\small
\begin{longtable}{@{}lll@{}}
\caption{The declared sample types and what the page makes of them.}\label{tab:d2-vocabulary} \\
\toprule \textbf{Type} & \textbf{Meaning in the specification} & \textbf{What the page reports} \\\midrule\endfirsthead
\toprule \textbf{Type} & \textbf{Meaning in the specification} & \textbf{What the page reports} \\\midrule\endhead
\texttt{unkn} & unknown sample & Unknown \\
\texttt{ntc} & non-template control & NTC \\
\texttt{nac} & no amplification control & not recognised: Unknown \\
\texttt{std} & standard sample & Standard \\
\texttt{ntp} & no target present & not recognised: Unknown \\
\texttt{nrt} & minus-RT control & not recognised: Unknown \\
\texttt{pos} & positive control & Positive control \\
\texttt{opt} & optical calibrator sample & reported as Unknown \\
\texttt{neg} & not defined by the specification & Negative control (unreachable) \\
\bottomrule
\end{longtable}}


A minus-RT control and a no-target-present control are blanks: nothing should
amplify in them. Reported as unknown samples they are not blanks to the review, so
a rising curve in one raises nothing --- which is the class of event this tool
exists to catch.

\vhead{Evidence}
A file carrying one reaction of each declared type, read by the page:

{\small
\begin{longtable}{@{}llllll@{}}
\toprule \textbf{Well} & \textbf{Sample} & \textbf{Declared} & \textbf{Role assigned} & \textbf{Value} & \textbf{Call} \\\midrule\endfirsthead
\toprule \textbf{Well} & \textbf{Sample} & \textbf{Declared} & \textbf{Role assigned} & \textbf{Value} & \textbf{Call} \\\midrule\endhead
\texttt{A01} & \texttt{S\_unkn} & \texttt{unkn} & Unknown & 20.5 & Analysed \\
\texttt{A02} & \texttt{S\_nac} & \texttt{nac} & Unknown & -1 & Analysed \\
\texttt{A03} & \texttt{S\_ntp} & \texttt{ntp} & Unknown & 33.2 & Analysed \\
\texttt{A04} & \texttt{S\_nrt} & \texttt{nrt} & Unknown & 31 & Analysed \\
\texttt{A05} & \texttt{S\_opt} & \texttt{opt} & Unknown & 18 & Analysed \\
\texttt{A06} & \texttt{S\_ntc} & \texttt{ntc} & NTC & 35.5 & Analysed \\
\bottomrule
\end{longtable}}

\vhead{Proposed correction}
\begin{lstlisting}[language=,numbers=left,firstnumber=1,basicstyle=\ttfamily\scriptsize\color{white}]
const RDML_SAMPLE_ROLE={unkn:"Unknown",ntc:"NTC",nac:"No-amplification control",
  std:"Standard",ntp:"No target present",nrt:"Minus-RT control",
  pos:"Positive control",opt:"Optical calibrator"};
\end{lstlisting}


\vhead{A correction to the earlier record}
The record delivered earlier stated that this mapping was complete against the
schema. It was checked against the corpus, which carries only two of the eight
types, and not against the schema. It is not complete.

\subsection*{D-3 --- the ``not available'' value is read as a measurement}
The specification records an absent result as $-1.0$. The reader converts the text
to a number and keeps it, so a reaction the file declares to have no result is
stored with a value of $-1$ and called analysed.

{\small
\begin{longtable}{@{}lrrrrrr@{}}
\caption{The not-available sentinel on the corpus.}\label{tab:d3-sentinel} \\
\toprule \textbf{File} & \textbf{Reactions} & \textbf{Below zero} & \textbf{Called analysed} & \textbf{Undetermined} & \textbf{N0 at -1} & \textbf{Efficiency at -1} \\\midrule\endfirsthead
\toprule \textbf{File} & \textbf{Reactions} & \textbf{Below zero} & \textbf{Called analysed} & \textbf{Undetermined} & \textbf{N0 at -1} & \textbf{Efficiency at -1} \\\midrule\endhead
\texttt{PAIRED-EXPE…} & 192 & 0 & 0 & 158 & 0 & 0 \\
\texttt{RELQUANT-PA…} & 72 & 0 & 0 & 6 & 0 & 0 \\
\texttt{example\_1\_r…} & 285 & 0 & 0 & 285 & 0 & 0 \\
\texttt{example\_2\_t…} & 285 & 0 & 0 & 285 & 0 & 0 \\
\texttt{example\_3\_l…} & 285 & 19 & 19 & 0 & 19 & 3 \\
\texttt{example\_4\_i…} & 285 & 19 & 19 & 0 & 19 & 3 \\
\texttt{example\_LCG…} & 50 & 0 & 0 & 50 & 0 & 0 \\
\bottomrule
\end{longtable}}


Two files are affected, 19 of 285 reactions in each. A value of $-1$ sorts below
every real one, so on a chart or in a trend it reads as the earliest crossing on
the plate. This is the same class of defect as the sentinel read as a crossing at
cycle zero, corrected on 23~September, and the property check that catches that one
does not catch this one. A further consequence: an efficiency of $-1$ is present in
three reactions of each file, and the cycle-ceiling rule is guarded by the absence
of an efficiency, so the guard is satisfied by a value that means ``not available''.

\vhead{Proposed correction}

\begin{lstlisting}[language=,numbers=left,firstnumber=1,basicstyle=\ttfamily\scriptsize\color{white}]
function rdmlNum(el,name){
  const t=rdmlText(el,name);if(t==="")return null;
  const n=Number(t);if(!Number.isFinite(n))return null;
  return n===-1?null:n;        /* -1.0 is the file saying "not available" */
}
\end{lstlisting}


This changes stored values --- 19 reactions per affected file move from analysed to
undetermined --- and therefore waits for agreement.

\subsection*{D-4 --- a sample declared per target keeps only its first declaration}
A sample may declare a different type for each target. The reader reads the first
declaration and applies it to every target of that sample.

{\small
\begin{longtable}{@{}llllll@{}}
\toprule \textbf{Well} & \textbf{Sample} & \textbf{Target} & \textbf{Declared for it} & \textbf{Role assigned} & \textbf{Value} \\\midrule\endfirsthead
\toprule \textbf{Well} & \textbf{Sample} & \textbf{Target} & \textbf{Declared for it} & \textbf{Role assigned} & \textbf{Value} \\\midrule\endhead
\texttt{A01} & \texttt{S\_dual} & \texttt{T1} & \texttt{std} & Standard & 22 \\
\texttt{A01} & \texttt{S\_dual} & \texttt{T2} & \texttt{ntp} & Standard & 34 \\
\bottomrule
\end{longtable}}


The no-target-present reaction is reported as a standard carrying a value of 34,
which would place it on a standard curve.

\subsection*{D-5 --- a free-format run loses every reaction}
The specification defines a row count of $-1$ to mean ``do not reconstruct a plate,
show the reactions as a list''. The reader multiplies rows by columns to bound the
plate, so the bound becomes negative and every reaction is discarded.

\begin{lstlisting}[language=,numbers=left,firstnumber=1,basicstyle=\ttfamily\scriptsize\color{white}]
declared rows -1, columns 1  ->  rows -1  columns 1  highest position -2  reactions kept 0
\end{lstlisting}


A run that declares no rows at all is silently given eight by the same expression.

\subsection*{D-6 --- rotor positions are labelled as plate wells}
A rotor declares one column and a numeric row label. The reader labels the
positions as if they were plate wells; the specification asks for plain numbers and
for no column label at all.

\subsection*{D-7 --- two accepted spellings of the format are rejected}
The format notes state that the archive may carry either of two extensions, and
that software should also read the uncompressed XML. The intake accepts one
extension and an archive that contains the expected entry; the other extension and
a bare XML file are rejected as unreadable. Both are one clause in the intake.

\subsection*{D-8 --- the file's own checksum is not used}
The format defines an identifier block carrying a publisher, a serial number and a
hash over the file. The reader reports, for every exchange file including one that
carries such a block:

\begin{lstlisting}[language=,numbers=left,firstnumber=1,basicstyle=\ttfamily\scriptsize\color{white}]
RDML carries no container checksum; the ZIP entry CRC-32 was verified on reading
\end{lstlisting}


That statement is wrong for a file that carries the block, and the page already
contains the hash implementation it would need. For a tool whose purpose is file
forensics this is the most valuable of the conformance findings.

\subsection*{D-9 --- the collection software is a structured element}
The element is defined as a name followed by a version. The reader takes the
element's text content, which concatenates both with the file's own indentation.
Two corpus files are affected. The correction is to read the two children and join
them with a space; it changes a displayed string only.

\subsection*{D-10 --- partitioned reactions are not read}
Partitioned reactions are recorded as counts, with the per-partition table stored
in a separate file inside the archive. The reader ignores the element, so such a
reaction contributes no result and nothing says why.

\subsection*{D-11 --- a latent key mismatch}
When a container and an exchange file are matched, the file-name key is normalised
but the experiment key is only lower-cased. No corpus file is affected --- the case
was tested by renaming an export so that only the experiment identifier could match,
and the match still held --- so this is recorded as a hardening, not a defect.

\section{D-12 --- the function index was short}
The index built by parsing the script recorded fewer symbols than the page exposes
at run time. The index in Appendix~\ref{app:index} is now built from the same walk
used for this part and records 676, including the members of the one module that
hides its helpers.

\section{Behaviour that is by design}
\vhead{Functions that change application state}

{\small
\begin{longtable}{@{}ll@{}}
\toprule \textbf{Function} & \textbf{What it changes} \\\midrule\endfirsthead
\toprule \textbf{Function} & \textbf{What it changes} \\\midrule\endhead
\texttt{appReadSuperseded} & global state changed during a read-only probe \\
\texttt{ccSetCfg} & global state changed during a read-only probe \\
\texttt{intake} & global state changed during a read-only probe \\
\texttt{mgReadFiles} & global state changed during a read-only probe \\
\texttt{readStaged} & global state changed during a read-only probe \\
\texttt{readStagedCore} & global state changed during a read-only probe \\
\texttt{sopSetPath} & global state changed during a read-only probe \\
\bottomrule
\end{longtable}}

\vhead{Entry points that are not called}
{\small
\begin{longtable}{@{}ll@{}}
\toprule \textbf{Symbol} & \textbf{Reason} \\\midrule\endfirsthead
\toprule \textbf{Symbol} & \textbf{Reason} \\\midrule\endhead
\texttt{boot} & entry point; re-runs the whole page \\
\texttt{ccImport} & replaces the instrument history by design \\
\texttt{invalidateAnalysisCaches} & clears caches by design \\
\texttt{setPhase} & mutates the lifecycle phase by design \\
\bottomrule
\end{longtable}}


\vhead{Helpers private to a module}
25 names are declared inside the decoder module's closure. They cannot be called by
name from the page; the module exports ten of its members, which are probed through
it, and all of them are printed in full with the module in
Chapter~\ref{ch:val-functions}.

\par{\small \texttt{CRC\_T}, \texttt{P}, \texttt{QS\_DEFAULT\_COLS}, \texttt{ROWS}, \texttt{argb}, \texttt{bracketList}, \texttt{fmt}, \texttt{hex}, \texttt{inflateRaw}, \texttt{ini}, \texttt{javaHashOrder}, \texttt{kid}, \texttt{kids}, \texttt{makeZip}, \texttt{pad2}, \texttt{parseLog}, \texttt{pill}, \texttt{qsCalDate}, \texttt{qsCols}, \texttt{qsHeader}, \texttt{qsRGB}, \texttt{qsStamp}, \texttt{txt}, \texttt{wellPos}, \texttt{xml}.}\par


\section{What was checked and found sound}
\begin{itemize}
\item Every property check holds, including the three that guard the defects
  corrected on 23~September.
\item Every check against the run whose answers were fixed first holds.
\item Every tab, review view, graph, console view and export builder runs without a
  page error, and the exports reproduce stored values unchanged.
\item The two readings of the one experiment present in both formats agree on every
  reaction they share, with no differing call.
\item Of 676 symbols, none raises an error on a stated input.
\end{itemize}

